% LaTeX master Datei(en) zusammengestellt von Falk-Jonatan Strube zur Nutzung an der Hochschule für Technik und Wirtschaft Dresden: https://github.com/genericFJS/htw
\documentclass{scrreprt}
\gdef\pathtomaster{../../_LaTeX_master}
%% HTW Dresden Corporate Design
%% Copyright (C) Falk-Jonatan Strube <falkjonatanstrube@gmail.com>, 2017
%%% === Anfang Paket htwscr ===
\makeatletter
\ifcsname KOMAoption\endcsname%
\KOMAoption{footheight}{11mm}	% mehr Platz für Wasserzeichen o.ä. in der Fußzeile
\recalctypearea	% s.o.
\RequirePackage{geometry}
\geometry{top=25mm, bottom=25mm, headsep=10mm, inner=25mm, outer=25mm, footskip=15mm, marginparsep=2mm}
\RequirePackage{helvet}	% Corporate Design schreibt Arial/Helvetica vor
\renewcommand{\familydefault}{\sfdefault}	% s.o.
\RequirePackage{graphicx}	% für Grafiken
\RequirePackage{xcolor}		% für Farben
\definecolor{htworange}{RGB}{249,155,28}	% Corporate Design orange CMYK(0/45/100/0)
\definecolor{htwblue}{RGB}{0,116,188}			% Corporate Design blau CMYK(90/50/0/0)
\definecolor{htwgrey}{RGB}{128,128,128}		% Corporate Design gray 60% Schwarz [Annahme: CMYK(0,0,0,60)]
\RequirePackage{ifthen,xifthen}	% für Bedingungen/Optionen
\RequirePackage{etoolbox}				% s.o.

%%% Format des HTW-Headers (wird hier als Texteelement verwenden) nach dem Style der PowerPoint-Vorlagen. Weiterhin inspiriert durch das Paket tudscr (Corporate Design der TU Dresden)
%%% Dateipfad des Logos abhängig von der Verwendung als Paket/Input
\ifdefined\@isHTWCDpackage%
	\gdef\@pathtoimage{HTW}%
\else%
	\ifdefined\pathtomaster%
		\gdef\@pathtoimage{\pathtomaster /HTW}%
	\else 
		\gdef\@pathtoimage{HTW}%
	\fi%
\fi%
\newcommand{\pathtologo}[1]{\gdef\@pathtoimage{#1}}
%%% Berechnung der Größe des Logos
\newlength{\htwLogoHeight}
\setlength{\htwLogoHeight}{12mm}	% Mindesthöhe des Logos laut Corporate Design: 11mm
\newlength{\@htwLogoWidth}
\AtBeginDocument{
\def\printLogo{\includegraphics[height=\htwLogoHeight]{\@pathtoimage}}
\setlength\@htwLogoWidth{\widthof{\printLogo}}
}
%%% Definition des Header
\newcommand{\@htwheader}{{\noindent%
%%% Header ist abhängig von dem Vorhandensein von department, institute und chair
\newcommand{\hasDepInsCh}{\gdef\@hasDepInsCh{}}%
\ifdefined\@department\hasDepInsCh\fi%
\ifdefined\@institute\hasDepInsCh\fi%
\ifdefined\@chair\hasDepInsCh\fi%
\begin{minipage}[b][\htwLogoHeight]{\dimexpr \textwidth - \@htwLogoWidth - 2mm \relax}% 2mm Abstand zwischen Logo und Linie
{\small{\bfseries%
\ifdefined\@faculty\@faculty{}\fi} %
\strut}%
{\color{htworange}\hrule height 1pt}%
\vspace{.1mm}%
\end{minipage}
\hfill%
\begin{minipage}[b][\htwLogoHeight]{\@htwLogoWidth}
\includegraphics[height=\htwLogoHeight]{\@pathtoimage}
\end{minipage}
\ifdefined\@hasDepInsCh%
\\%
\begin{minipage}{\textwidth}
\small%
\ifdefined\@department\@department{}\fi%
\ifdefined\@institute\ifdefined\@department,~\fi \@institute\fi%
\ifdefined\@chair\ifdefined\@department,~\else\ifdefined\@institute,~\fi\fi\@chair\fi\strut
%{\color{htworange}\hrule height 0.5pt}
\end{minipage}
\fi
}}

%%% Definition der Fußzeilen
\usepackage{scrlayer-scrpage}
\@ifclassloaded{scrreprt}{%
\renewcommand*{\partpagestyle}{scrheadings}
\renewcommand*{\chapterpagestyle}{scrheadings}%
}{}
\pagestyle{scrheadings}
\cfoot{\vfill\pagemark}
\chead{}

%%% Einstellungen
\newcommand{\simpledate}{\gdef\@simpledate{}}		% Datum nicht mit "Eingereicht am:" Vorsatz
\newcommand{\nouppercase}{\gdef\@nouppercase{}}	% Titel/Parts/Chapters usw. nicht in Großbuchstaben

%%% Einiges erst zu Beginn des Dokumentes ausführen, um Einstellungen zu berücksichtigen
\AtBeginDocument{
\renewcommand\SS{SS}	% fix big ß
%%% Sektionen usw. ggf. in Großbuchstaben:
\providecommand{\chapterlinesformat}[3]{}	% compatibility
\providecommand{\chapterlineswithprefixformat}[3]{}	% compatibility
\providecommand{\sectionlinesformat}[4]{}	% compatibility
\ifdefined\@nouppercase\relax\else
\renewcommand\sectionlinesformat[4]{\@hangfrom{\hskip#2 #3}{\MakeUppercase{#4}}}
\@ifclassloaded{scrreprt}{%
\renewcommand\chapterlinesformat[3]{\@hangfrom{#2}{\MakeUppercase{#3}}}
\renewcommand\chapterlineswithprefixformat[3]{\MakeUppercase{#2#3}}%
}{}
\fi
%%% Anpassung der Part-Seite
\renewcommand*{\partheadstartvskip}{\@htwheader\vspace*{30mm}}
\renewcommand*{\raggedpart}{\raggedright}
\renewcommand*{\partformat}{\huge\ifdefined\@nouppercase \partname~\thepart\else\MakeUppercase{\partname~\thepart}\fi.}
\ifdefined\@nouppercase\setkomafont{part}{\Huge}\else\setkomafont{part}{\Huge\MakeUppercase}\fi
}

\newcommand{\bilingual}[3]{\ifdefined #1 \let #1\undefined\fi\newcommand{#1}{\iflanguage{english}{#3}{#2}}}	% Übersetzung von einigen Variablen in mehrere Sprachen

\newcommand{\checkThesis}[2]{\ifthenelse{\equal{#1}{diss}\OR\equal{#1}{doctoral}\OR\equal{#1}{phd}}{
\gdef#2{\dissertationname}}{
\ifthenelse{\equal{#1}{diploma}}{
\gdef#2{\diplomathesisname}}{
\ifthenelse{\equal{#1}{master}}{
\gdef#2{\masterthesisname}}{
\ifthenelse{\equal{#1}{bachelor}}{
\gdef#2{\bachelorthesisname}}{
\ifthenelse{\equal{#1}{student}}{
\gdef#2{\studentthesisname}}{
\ifthenelse{\equal{#1}{evidence}}{
\gdef#2{\studentresearchname}}{
\ifthenelse{\equal{#1}{project}}{
\gdef#2{\projectpapername}}{
\ifthenelse{\equal{#1}{seminar}}{
\gdef#2{\seminarpapername}}{
\ifthenelse{\equal{#1}{term}}{
\gdef#2{\termpapername}}{
\ifthenelse{\equal{#1}{research}}{
\gdef#2{\researchname}}{
\ifthenelse{\equal{#1}{log}}{
\gdef#2{\logname}}{
\ifthenelse{\equal{#1}{report}}{
\gdef#2{\reportname}}{
\ifthenelse{\equal{#1}{internship}}{
\gdef#2{\internshipname}}{
\ifthenelse{\equal{#1}{lecture}\OR\equal{#1}{lesson}\OR\equal{#1}{pract}}{
\bilingual{\professorname}{Vorlesung von}{Lecture by}
\simpledate
\gdef\@islecture{\relax}
\ifthenelse{\equal{#1}{lecture}}{\gdef#2{\lecturename}}{
\ifthenelse{\equal{#1}{lesson}}{\gdef#2{\lessonname}}{
\gdef#2{\practname}
}}}{
\gdef#2{#1}}}}}}}}}}}}}}}}

%%% Textfragmente der Titelseite in deutsch und englisch
\bilingual{\dateofbirthname}{Geboren am:}{Born on:}
\bilingual{\placeofbirthname}{in}{in}
\bilingual{\coursename}{Studiengang:}{Course:}
\bilingual{\disciplinename}{Studienrichtung:}{Discipline:}
\bilingual{\matriculationnumbername}{Matrikelnummer:}{Matriculation number:}
\bilingual{\matriculationyearname}{Immatrikulationsjahr}{Matriculation year:}
\bilingual{\graduationname}{zur Erlangung des akademischen Grades}{to achieve the academic degree}
\bilingual{\refereename}{Gutachter}{Referee}
\bilingual{\advisorname}{Fachreferent}{Advisor}
\bilingual{\supervisorname}{Betreuer}{Supervisor}
\bilingual{\professorname}{Betreuender Hochschullehrer}{Supervising professor}
\bilingual{\datename}{Eingereicht am:}{Submitted on:}
%% Thesis Typen (inspiriert von dem Paket tudscr):
\bilingual{\dissertationname}{Dissertation}{Dissertation} % diss, doctoral, phd 
\bilingual{\diplomathesisname}{Diplomarbeit}{Diploma Thesis} % diploma 
\bilingual{\masterthesisname}{Master-Arbeit}{Master Thesis} % master 
\bilingual{\bachelorthesisname}{Bachelor-Arbeit}{Bachelor Thesis} % bachelor 
\bilingual{\studentthesisname}{Studienarbeit}{Student Thesis} % student 
\bilingual{\studentresearchname}{Großer Beleg}{Student Research Project} % evidence 
\bilingual{\projectpapername}{Projektarbeit}{Project Paper} % project
\bilingual{\seminarpapername}{Seminararbeit}{Seminar Paper} % seminar
\bilingual{\termpapername}{Hausarbeit}{Term Paper} % term
\bilingual{\researchname}{Forschungsbericht}{Research Report} % research 
\bilingual{\logname}{Protokoll}{Log} % log 
\bilingual{\reportname}{Bericht}{Report} % report
\bilingual{\internshipname}{Praktikumsbericht}{Internship Report} % internship 
\bilingual{\lecturename}{Vorlesungsmitschrift}{Lecture Notes} % lecture 
\bilingual{\lectureauthorname}{Mitschrift von}{Notes by} % lecture  
\bilingual{\lessonname}{\"Ubungsmitschrift}{Lesson Notes} % lecture  
\bilingual{\practname}{Praktikumsmitschrift}{Pracital Course Notes} % lecture 
%%% Variablen für die Titelseite (inspiriert von dem Paket tudscr):
\renewcommand{\subject}[1]{\checkThesis{#1}{\@subject}}
\newcommand{\faculty}[1]{\gdef\@faculty{#1}}
\newcommand{\department}[1]{\gdef\@department{#1}}
\newcommand{\institute}[1]{\gdef\@institute{#1}}
\newcommand{\chair}[1]{\gdef\@chair{#1}}
\newcommand{\authormore}[1]{\gdef\@authormore{#1}}
\newcommand{\dateofbirth}[1]{\gdef\@dateofbirth{#1}}
\newcommand{\placeofbirth}[1]{\gdef\@placeofbirth{#1}}
\newcommand{\course}[1]{\gdef\@course{#1}}
\newcommand{\discipline}[1]{\gdef\@discipline{#1}}
\newcommand{\matriculationnumber}[1]{\gdef\@matriculationnumber{#1}}
\newcommand{\matriculationyear}[1]{\gdef\@matriculationyear{#1}}
\newcommand{\thesis}[1]{\checkThesis{#1}{\@thesis}}
\newcommand{\graduation}[2][]{\gdef\@graduation{#2}\gdef\@graduationshort{#1}}
\newcommand{\supervisor}[1]{\gdef\@supervisor{#1}}
\newcommand{\referee}[1]{\gdef\@referee{#1}}
\newcommand{\advisor}[1]{\gdef\@advisor{#1}}
\newcommand{\professor}[1]{\gdef\@professor{#1}}
%%% Titelseite (inspiriert von dem Paket tudscr):
\renewcommand*{\maketitle}{%
\begin{titlepage}%
\begin{flushleft}
\renewcommand{\thefootnote}{$\star$}%
%%% Kopfzeile HTW-Style:
\@htwheader\\[13mm]
%%% Die Zeilen für titlehead, subject, title, subtitle und author sind immer vorhanden (spacing)
{\noindent\@titlehead{}\strut}\\[3mm]%
{\Large\@subject{}\strut}\\[8mm]%
{\Huge\bfseries%
\ifdefined\@nouppercase\@title{}\else\MakeUppercase{\@title{}}\fi%
\strut}\\
{\Large\bfseries%
\ifdefined\@nouppercase\@subtitle{}\else\MakeUppercase{\@subtitle{}}\fi%
\strut}\\[11mm]
\ifdefined\@islecture\strut\else{\Large\@author{}\strut}\fi % Autor wird bei Vorlesungsmitschrift unten angezeigt
%%% Weitere Felder werden nur angezeigt, wenn sie definiert wurden
\ifdefined\@authormore\ifdefined\@islecture\else\\\fi{\@authormore{}\strut}\fi
\ifdefined\@dateofbirth\\\dateofbirthname{} \@dateofbirth{} \ifdefined\@placeofbirth\placeofbirthname{} \@placeofbirth{}\strut\fi\fi
\ifdefined\@course\\\coursename{} \@course{}\strut\fi
\ifdefined\@discipline\\\disciplinename{} \@discipline{}\strut\fi
\ifdefined\@matriculationnumber\\\matriculationnumbername{} \@matriculationnumber{}\strut\fi
\ifdefined\@matriculationyear\\\matriculationyearname{} \@matriculationyear{}\strut\fi
\ifdefined\@thesis\\[10mm]{\LARGE\bfseries
\ifdefined\@nouppercase\@thesis{}\else\MakeUppercase{\@thesis{}}\fi
\strut}
\ifdefined\@graduation\\[2mm]{\graduationname{}\strut}
\\[5mm]{\LARGE\bfseries
\ifdefined\@nouppercase\relax\else\MakeUppercase{\fi%
\@graduation{} \ifdefined\@graduationshort(\@graduationshort{})\fi
\ifdefined\@nouppercase\relax\else}\fi\strut}\fi\fi
\vfill\noindent
\ifdefined\@islecture\\[1mm]{\lectureauthorname{}\strut}\\[.7mm]{\Large\@author{}\strut}\fi
\ifdefined\@referee{\refereename{}\strut}\\[.7mm]{\Large\@referee{}\strut}\fi
\ifdefined\@advisor\\[1mm]{\advisorname{}\strut}\\[.7mm]{\Large\@advisor{}\strut}\fi
\ifdefined\@supervisor\\[1mm]{\supervisorname{}\strut}\\[.7mm]{\Large\@supervisor{}\strut}\fi
\ifdefined\@professor\\[1mm]{\professorname{}\strut}\\[.7mm]{\Large\@professor{}\strut}\fi
\ifdefempty{\@date}{}{\ifdefined\@simpledate\\[6mm]\relax\else\\[3mm]\datename{}~\fi\@date{}\strut}
\ifdefempty{\@publishers}{}{\\[4.5mm]{\Large\@publishers{}\strut}}%
\renewcommand\footnoterule{}
\@thanks\global\let\@thanks\@empty
\end{flushleft}
\end{titlepage}
\setcounter{footnote}{0}
}
\else%
\PackageError{htwcd}{htwcd needs a KOMA class!}%
\fi%
\makeatother
%%% === Ende Paket htwscr ===
%-- Nummerierung bis Subsubsection für Report
\setcounter{secnumdepth}{3}	% zählt auch subsubsection
\setcounter{tocdepth}{3}	% Inhaltsverzeichnis bis in subsubsection

% LuaLaTeX Unterscheidung
\usepackage{ifluatex}

% Input inkl. Umlaute, Silbentrennung
\ifluatex
	\usepackage{fontspec}
\else
	\usepackage[utf8]{inputenc}
	\usepackage[T1]{fontenc}
\fi
\usepackage[english,ngerman]{babel}
\usepackage{csquotes}	% Anführungszeichen
\usepackage{marvosym}
\usepackage{eurosym}

% Style-Aufhübschung
\usepackage{soul, color}	% Kapitälchen, Unterstrichen, Durchgestrichen usw. im Text
%\usepackage{titleref}
\usepackage[breakwords, fit]{truncate}	% Abschneiden von Sätzen
\renewcommand{\TruncateMarker}{\,…}

% Mathe usw.
\usepackage{amssymb}
\usepackage{amsthm}
\usepackage[fleqn,intlimits]{amsmath}	% fleqn: align-Umgebung rechtsbündig; intlimits: Integralgrenzen immer ober-/unterhalb
%\usepackage{mathtools} % u.a. schönere underbraces
\usepackage{xcolor}
\usepackage{esint}	% Schönere Integrale, \oiint vorhanden
\everymath=\expandafter{\the\everymath\displaystyle}	% Mathe Inhalte werden weniger verkleinert
\usepackage{wasysym}	% mehr Symbole, bspw \lightning
%\renewcommand{\int}{\int\limits}
%\usepackage{xfrac}	% mehr fracs: sfrac{}{}
\let\oldemptyset\emptyset	% schöneres emptyset
\let\emptyset\varnothing
%\RequirePackage{mathabx}	% mehr Symbole
\mathchardef\mhyphen="2D	% Hyphen in Math
\usepackage{mathtools}		% zum Beispiel {pmatrix}[r] für Formatierung der Matrix

% tikz usw.
\usepackage{tikz}
\usepackage{pgfplots}
\pgfplotsset{compat=1.11}	% Umgeht Fehlermeldung
\usetikzlibrary{graphs}
%\usetikzlibrary{through}	% ???
\usetikzlibrary{arrows}
\usetikzlibrary{arrows.meta}	% Pfeile verändern / vergrößern: \draw[-{>[scale=1.5]}] (-3,5) -> (-3,3);
\usetikzlibrary{automata,positioning} % Zeilenumbruch im Node node[align=center] {Text\\nächste Zeile} automata für Graphen
\usetikzlibrary{matrix}
\usetikzlibrary{patterns}	% Schraffierte Füllung
\usetikzlibrary{shapes.geometric}	% Polygon usw.
\tikzstyle{reverseclip}=[insert path={	% Inverser Clip \clip
	(current page.north east) --
	(current page.south east) --
	(current page.south west) --
	(current page.north west) --
	(current page.north east)}
% Nutzen: 
%\begin{tikzpicture}[remember picture]
%\begin{scope}
%\begin{pgfinterruptboundingbox}
%\draw [clip] DIE FLÄCHE, IN DER OBJEKT NICHT ERSCHEINEN SOLL [reverseclip];
%\end{pgfinterruptboundingbox}
%\draw DAS OBJEKT;
%\end{scope}
%\end{tikzpicture}
]	% Achtung: dafür muss doppelt kompliert werden!
\usepackage{graphpap}	% Grid für Graphen
\tikzset{every state/.style={inner sep=2pt, minimum size=2em}}
\usetikzlibrary{mindmap, backgrounds}
%\usepackage{tikz-uml}	% braucht Dateien: http://perso.ensta-paristech.fr/~kielbasi/tikzuml/

% Tabular
\usepackage{longtable}	% Große Tabellen über mehrere Seiten
\usepackage{multirow}	% Multirow/-column: \multirow{2[Anzahl der Zeilen]}{*[Format]}{Test[Inhalt]} oder \multicolumn{7[Anzahl der Reihen]}{|c|[Format]}{Test2[Inhalt]}
\renewcommand{\arraystretch}{1.3} % Tabellenlinien nicht zu dicht
\usepackage{colortbl}
\arrayrulecolor{gray}	% heller Tabellenlinien
\usepackage{array}	% für folgende 3 Zeilen (für Spalten fester breite mit entsprechender Ausrichtung):
\newcolumntype{L}[1]{>{\raggedright\let\newline\\\arraybackslash\hspace{0pt}}m{\dimexpr#1\columnwidth-2\tabcolsep-1.5\arrayrulewidth}}
\newcolumntype{C}[1]{>{\centering\let\newline\\\arraybackslash\hspace{0pt}}m{\dimexpr#1\columnwidth-2\tabcolsep-1.5\arrayrulewidth}}
\newcolumntype{R}[1]{>{\raggedleft\let\newline\\\arraybackslash\hspace{0pt}}m{\dimexpr#1\columnwidth-2\tabcolsep-1.5\arrayrulewidth}}
\usepackage{booktabs}	% Kompatibilität für array und \hrule
\usepackage{caption}	% Um auch unbeschriftete Captions mit \caption* zu machen

% Nützliches
\usepackage{verbatim}	% u.a. zum auskommentieren via \begin{comment} \end{comment}
\usepackage{tabto}	% Tabs: /tab zum nächsten Tab oder /tabto{.5 \CurrentLineWidth} zur Stelle in der Linie
\NumTabs{6}	% Anzahl von Tabs pro Zeile zum springen
\usepackage{listings} % Source-Code mit Tabs
\usepackage{lstautogobble} 
\usepackage{enumitem}	% Anpassung der enumerates
%\setlist[enumerate,1]{label=(\arabic*)}	% global andere Enum-Items
\renewcommand{\labelitemiii}{$\scriptscriptstyle ^\blacklozenge$} % global andere 3. Item-Aufzählungszeichen
\usepackage{letltxmacro} % neue Definiton von Grundbefehlen
% Nutzen:
%\LetLtxMacro{\oldemph}{\emph}
%\renewcommand{\emph}[1]{\oldemph{#1}}
\RequirePackage{xpatch}	% ua. Konkatenieren von Strings/Variablen (etoolbox)
\usepackage{xstring}	% String Operationen
\usepackage{minibox}	% Minibox anstatt \fbox{} für Boxen mit (manuellen) Zeilenumbruch
\usepackage{framed}		% \begin{framed} statt minibox/fbox mit automatischem Zeilenumbruch
\usepackage{microtype}	% Verfeinerte Hyphenation
\usepackage{import}		
\usepackage{standalone}	% Zusammen mit import: Einfügen von separaten Latex-Dateien (die schon header haben)

% Einrichtung von lst
%\renewcommand*{\ttdefault}{qcr}
\lstset{
basicstyle=\ttfamily, 
%basicstyle=\footnotesize\sffamily\color{black},
%mathescape=true, 
%escapeinside=^^, 
autogobble, 
tabsize=2,
frame=single,
rulecolor=\color{lightgray},
numbers=left,
numbersep=5pt,
numberstyle=\tiny\color{gray},
commentstyle=\color{gray},
keywordstyle=\color{htwblue},
stringstyle=\color{htworange},
morecomment=[l][\color{magenta}]{\#}
showspaces=false,
showstringspaces=false,
breaklines=true,
literate=%
    {Ö}{{\"O}}1
    {Ä}{{\"A}}1
    {Ü}{{\"U}}1
    {ß}{{\ss}}1
    {ü}{{\"u}}1
    {ä}{{\"a}}1
    {ö}{{\"o}}1
    {„}{{\glqq}}1
    {“}{{\grqq}}1
    {~}{{\textasciitilde}}1
}
\usepackage{scrhack} % Fehler umgehen
\def\ContinueLineNumber{\lstset{firstnumber=last}} % vor lstlisting. Zum wechsel zum nicht-kontinuierlichen muss wieder \StartLineAt1 eingegeben werden
\def\StartLineAt#1{\lstset{firstnumber=#1}} % vor lstlisting \StartLineAt30 eingeben, um bei Zeile 30 zu starten
\let\numberLineAt\StartLineAt

% BibTeX
\usepackage[
backend=bibtex,
bibencoding=ascii,
style=authoryear, citestyle=authoryear-ibid,
%doi=false,
%isbn=false,
%url=false
]{biblatex}	% BibTeX
\usepackage{makeidx}
%\makeglossary
%\makeindex
\usepackage{url}	% Darstellung von URLs

% Grafiken
\usepackage{graphicx}
\usepackage{epstopdf}	% eps-Vektorgrafiken einfügen
\usepackage{transparent}	% transparent nutzen: {\transparent{0.4} ...}
%\epstopdfsetup{outdir=\customDir}
% Prüft, ob Grafik existiert (mit \ifvalidimage{}{}) [Quelle: https://tex.stackexchange.com/a/99176]:
\makeatletter
\newif\ifgraphicexist
\catcode`\*=11
\newcommand\ifvalidimage[1]{%
    \begingroup
    \global\graphicexisttrue
    \let\input@path\Ginput@path
    \filename@parse{#1}%
    \ifx\filename@ext\relax
    \@for\Gin@temp:=\Gin@extensions\do{%
        \ifx\Gin@ext\relax
        \Gin@getbase\Gin@temp
        \fi}%
    \else
    \Gin@getbase{\Gin@sepdefault\filename@ext}%
    \ifx\Gin@ext\relax
    \global\graphicexistfalse
    \def\Gin@base{\filename@area\filename@base}%
    \edef\Gin@ext{\Gin@sepdefault\filename@ext}%
    \fi
    \fi
    \ifx\Gin@ext\relax
    \global\graphicexistfalse
    \else 
    \@ifundefined{Gin@rule@\Gin@ext}%
    {\global\graphicexistfalse}%
    {}%
    \fi  
    \ifx\Gin@ext\relax 
    \gdef\imageextension{unknown}%
    \else
    \xdef\imageextension{\Gin@ext}%
    \fi 
    \endgroup 
    \ifgraphicexist
    \expandafter \@firstoftwo
    \else
    \expandafter \@secondoftwo
    \fi 
} 
\catcode`\*=12
\makeatother
\LetLtxMacro{\forceincludegraphics}{\includegraphics}	% neuer Befehl für includegraphics
\renewcommand{\includegraphics}[2][]{	% altes includegraphics neu definieren, damit es auch nicht vorhandene einfügt
\ifvalidimage{#2}{
\forceincludegraphics[#1]{#2}
}{
\message{Achtung: Grafik wurde nicht gefunden: '#2'}
\minibox[frame]{
\textbf{\StrSubstitute{#2}{_}{\_}}%  \ifthenelse{\isempty{#1}}{}{\\\textit{#1}}
}}}

% pdf-Setup
\usepackage{pdfpages}
\usepackage[bookmarks,%
bookmarksopen=false,% Klappt die Bookmarks in Acrobat aus
colorlinks=true,%
linkcolor=black,%
citecolor=black,%
urlcolor=htworange,%
]{hyperref}

% Nach Dokumentenbeginn ausführen:
\makeatletter
\AtBeginDocument{	
	% pdf Informationen
	\hypersetup{
		pdfauthor={\@author},
		pdftitle={\@title},
		pdfproducer={LaTeX Layout erstellt von Falk-Jonatan Strube}
	}
	% signierte Fußzeile
	%\providecommand{\ofootname}{Mitschrift von}
	\providecommand{\pathtomaster}{../}
	\ofoot{%
	%\ifdefined\ofoottext
	%\else
	%\begin{minipage}{40mm}
	%\begin{flushright}
	%\textcolor{gray}{
	%\ofootname\\
	%Falk-Jonatan Strube\vspace*{-1mm}}
	%\end{flushright}
	%\end{minipage}
	%\fi%
	%\hspace*{1mm}%
	\begin{minipage}{10 mm}
	\includegraphics[height=9mm]{\pathtomaster/FJS-logo_gray}
	\end{minipage}%
	}
}
\makeatother



%%% ==================================================
%%% ===================== MACROS =====================
%%% ==================================================
% Horizontale Linie:
\newcommand{\HRule}[1][\medskipamount]{\par
  \vspace*{\dimexpr-\parskip-\baselineskip+#1}
  \noindent\rule[0.2ex]{\linewidth}{0.2mm}\par
  \vspace*{\dimexpr-\parskip-.5\baselineskip+#1}}
% Gestrichelte horizontale Linie:
\RequirePackage{dashrule}
\newcommand{\HDRule}[1][\medskipamount]{\par
  \vspace*{\dimexpr-\parskip-\baselineskip+#1}
  \noindent\hdashrule[0.2ex]{\linewidth}{0.2mm}{1mm} \par
  \vspace*{\dimexpr-\parskip-.5\baselineskip+#1}}
% Mathe in Anführungszeichen:
\newsavebox{\mathbox}\newsavebox{\mathquote}
\makeatletter
\newcommand{\mq}[1]{% \mathquotes{<stuff>}
  \savebox{\mathquote}{\text{"}}% Save quotes
  \savebox{\mathbox}{$\displaystyle #1$}% Save <stuff>
  \raisebox{\dimexpr\ht\mathbox-\ht\mathquote\relax}{"}#1\raisebox{\dimexpr\ht\mathbox-\ht\mathquote\relax}{''}
}
\makeatother

% Paragraph mit Zähler (Section-Weise)
\newcounter{cparagraphC}
\newcommand{\cparagraph}[1]{
\stepcounter{cparagraphC}
\paragraph{\thesection{}-\thecparagraphC{} #1}
%\addcontentsline{toc}{subsubsection}{\thesection{}-\thecparagraphC{} #1}
\label{\thesection-\thecparagraphC}
}
\makeatletter
\@addtoreset{cparagraphC}{section}
\makeatother


% (Vorlesungs-)Folien einbinden:
\makeatletter
% Folien-Pfad (inkl. Dateiname ohne Endung und ggf. ohne Nummerierung)
\newcommand{\@slidePath}{}
\newcommand{\slidePath}[1]{\renewcommand{\@slidePath}{#1}}
% Folien Eigenschaften
\newcommand{\@slideScale}{0.5}
\newcommand{\slideScale}[1]{\renewcommand{\@slideScale}{#1}}
\AtBeginDocument{
% Folien von einer Datei skaliert
\newcommand{\slide}[2][\@slideScale]{\slides[#1]{}{#2}}
\newcommand{\slideTrim}[6][\@slideScale]{\slides[#1 , clip,  trim = #5cm #4cm #6cm #3cm]{}{#2}}
% Folien von mehreren nummerierten Dateien skaliert
\newcommand{\slides}[3][\@slideScale]{\begin{center}
\includegraphics[page=#3, scale=#1]{\@slidePath #2.pdf}
\end{center}}
}
\makeatother

% \emph{} anders definieren
\makeatletter
\DeclareRobustCommand{\em}{%
  \@nomath\em \if b\expandafter\@car\f@series\@nil
  \normalfont \else \scshape \fi}
\makeatother

% unwichtiges
\newcommand{\unimptnt}[1]{{\transparent{0.5}#1}}

% alph. enumerate
\newenvironment{anumerate}{\begin{enumerate}[label=(\alph*)]}{\end{enumerate}} % Alphabetische Aufzählung

% Hanging parameters
\newcommand{\hangpara}[1]{\par\noindent\hangindent+2em\hangafter=1 #1\par\noindent}

%% EINFACHE BEFEHLE

% Abkürzungen Mathe
\newcommand{\EE}{\mathbb{E}}
\newcommand{\QQ}{\mathbb{Q}}
\newcommand{\RR}{\mathbb{R}}
\newcommand{\CC}{\mathbb{C}}
\newcommand{\NN}{\mathbb{N}}
\newcommand{\ZZ}{\mathbb{Z}}
\newcommand{\PP}{\mathbb{P}}
\renewcommand{\SS}{\mathbb{S}}
\newcommand{\cA}{\mathcal{A}}
\newcommand{\cB}{\mathcal{B}}
\newcommand{\cC}{\mathcal{C}}
\newcommand{\cD}{\mathcal{D}}
\newcommand{\cE}{\mathcal{E}}
\newcommand{\cF}{\mathcal{F}}
\newcommand{\cG}{\mathcal{G}}
\newcommand{\cH}{\mathcal{H}}
\newcommand{\cI}{\mathcal{I}}
\newcommand{\cJ}{\mathcal{J}}
\newcommand{\cM}{\mathcal{M}}
\newcommand{\cN}{\mathcal{N}}
\newcommand{\cP}{\mathcal{P}}
\newcommand{\cR}{\mathcal{R}}
\newcommand{\cS}{\mathcal{S}}
\newcommand{\cZ}{\mathcal{Z}}
\newcommand{\cL}{\mathcal{L}}
\newcommand{\cT}{\mathcal{T}}
\newcommand{\cU}{\mathcal{U}}
\newcommand{\cX}{\mathcal{X}}
\newcommand{\cV}{\mathcal{V}}
\renewcommand{\phi}{\varphi}
\renewcommand{\epsilon}{\varepsilon}
\renewcommand{\theta}{\vartheta}

% Verschiedene als Mathe-Operatoren
\DeclareMathOperator{\arccot}{arccot}
\DeclareMathOperator{\arccosh}{arccosh}
\DeclareMathOperator{\arcsinh}{arcsinh}
\DeclareMathOperator{\arctanh}{arctanh}
\DeclareMathOperator{\arccoth}{arccoth} 
\DeclareMathOperator{\var}{Var} % Varianz 
\DeclareMathOperator{\cov}{Cov} % Co-Varianz 

% Farbdefinitionen
\definecolor{red}{RGB}{180,0,0}
\definecolor{green}{RGB}{75,160,0}
\definecolor{blue}{RGB}{0,75,200}
\definecolor{orange}{RGB}{255,128,0}
\definecolor{yellow}{RGB}{255,245,0}
\definecolor{purple}{RGB}{75,0,160}
\definecolor{cyan}{RGB}{0,160,160}
\definecolor{brown}{RGB}{120,60,10}

\definecolor{itteny}{RGB}{244,229,0}
\definecolor{ittenyo}{RGB}{253,198,11}
\definecolor{itteno}{RGB}{241,142,28}
\definecolor{ittenor}{RGB}{234,98,31}
\definecolor{ittenr}{RGB}{227,35,34}
\definecolor{ittenrp}{RGB}{196,3,125}
\definecolor{ittenp}{RGB}{109,57,139}
\definecolor{ittenpb}{RGB}{68,78,153}
\definecolor{ittenb}{RGB}{42,113,176}
\definecolor{ittenbg}{RGB}{6,150,187}
\definecolor{itteng}{RGB}{0,142,91}
\definecolor{ittengy}{RGB}{140,187,38}

% Textfarbe ändern
\newcommand{\tred}[1]{\textcolor{red}{#1}}
\newcommand{\tgreen}[1]{\textcolor{green}{#1}}
\newcommand{\tblue}[1]{\textcolor{blue}{#1}}
\newcommand{\torange}[1]{\textcolor{orange}{#1}}
\newcommand{\tyellow}[1]{\textcolor{yellow}{#1}}
\newcommand{\tpurple}[1]{\textcolor{purple}{#1}}
\newcommand{\tcyan}[1]{\textcolor{cyan}{#1}}
\newcommand{\tbrown}[1]{\textcolor{brown}{#1}}

% Umstellen der Tabellen Definition
\newcommand{\mpb}[1][.3]{\begin{minipage}{#1\textwidth}\vspace*{3pt}}
\newcommand{\mpe}{\vspace*{3pt}\end{minipage}}

\newcommand{\resultul}[1]{\underline{\underline{#1}}}
\newcommand{\parskp}{$ $\\}	% new line after paragraph
\newcommand{\parnl}{\unskip\parfillskip 0pt \par} % justified new line (brace line!)

% mathematische Abkürzungen
\newcommand{\corr}{\;\widehat{=}\;}
\newcommand{\mdeg}{^{\circ}}
\newcommand{\nok}[2]{\binom{#1}{#2}}	% n über k BESSER: \binom{n}{k}
\newcommand{\mtr}[1]{\begin{pmatrix}#1\end{pmatrix}}	% Matrix
\newcommand{\dtr}[1]{\begin{vmatrix}#1\end{vmatrix}}	% Determinante (Betragsmatrix)
\LetLtxMacro{\originalVec}{\vec}
\renewcommand{\vec}[1]{\underline{#1}}	% Vektorschreibweise
\newcommand{\imptnt}[1]{\colorbox{red!30}{#1}}	% Wichtiges
\newcommand{\intd}[1]{\,\mathrm{d}#1}
\newcommand{\diffd}[1]{\mathrm{d}#1}
% für Module-Rechnung: \pmod{}
\newcommand{\unit}[1]{\,\mathrm{#1}}
%\LetLtxMacro{\ntilde}{\tilde}
%\renewcommand{\tilde}{\widetilde}

% weiter Abkürzungen
\newcommand{\gdw}{genau dann wenn}
\newcommand{\lecdate}[1]{\marginpar{\small\textcolor{gray}{Vorlesung\newline #1}}}

\faculty{Fakultät Informatik/Mathematik}
\chair{Lehrstuhl für Künstliche Intelligenz / Theoretische Informatik}
\subject{Praktikumsmitschrift}
\title{Künstliche Intelligenz}
\author{Falk-Jonatan Strube}
\professor{Prof. Dr. Boris Hollas}
\renewcommand{\professorname}{Übung von}
\simpledate
\makeatletter
\gdef\@islecture{}
\makeatother

\begin{document}

\maketitle
\tableofcontents

\chapter{Praktikum 1: Prädikatenlogik}

\section{Aufgabe 1}
\begin{anumerate}
\item\label{1-1-a} Für alle $x$ gilt: Wenn $x$ eine Anfrage ist, gibt es ein $y$ in der entsprechenden Antwort.\\
Formel ist \emph{wahr}.
\item Für alle $x$ existiert ein $y$ für das gilt: Wenn $x$ eine Anfrage ist, gibt es eine entsprechende Antwort mit $x$ und $y$ (ist äquivalent mit \ref{1-1-a}).\\
Formel ist \emph{wahr}.
\item Es existiert ein $y$ bei dem für alle $x$ gilt: Wenn $x$ eine Anfrage ist, gibt es eine entsprechende Antwort mit $x$ und $y$ (Es gibt ein $y$, das für alle Anfragen die richtige Antwort ist).\\
Formel ist \emph{falsch}.\\
Beweis:\\
Angenommen, es gibt ein $y$, so dass $\forall x\,(anfrage(x) \to antwort(x,y))$ wahr ist. Wegen $(abfahrt, 9:00) \in antwort$ muss dann $y=9:00$ sein. Wegen $(ankunft, 12:00)\in antwort$ muss aber auch $y=12:00$. Widerspruch!
\end{anumerate}

\section{Aufgabe 2}
\begin{anumerate}
\item 
\begin{itemize}
\item $F_1$ ist \emph{wahr}:\\
Für alle Menschen und Fächer gilt: Wenn es die Kombination Mensch-Fach gibt, ist dieser Mensch Student.
\item $F_2$ ist \emph{falsch}:\\
Für alle Menschen und Fächer gilt: Die Kombination Mensch-Fach existiert und der Mensch ist Student.
\item $F_3$ ist \emph{wahr}:\\
Für alle Menschen existiert ein $y$ für das gilt: Wenn es die Kombination Mensch-$y$(bspw. auch „Toastbrot“) gibt, ist dieser Mensch Student.
\end{itemize}
\item 
\begin{itemize}
\item $F_1$ ist \emph{falsch}.
\item $F_2$ ist \emph{falsch}.
\item $F_3$ ist \emph{wahr}:\\
Es existiert ein ein Fach für jeden Menschen, das er nicht studiert.
\end{itemize}
\end{anumerate}
\begin{itemize}
\item $F_2$ bedeutet nicht „Wer ein Fach studiert, ist ein Student“, weil man die Formel umstellen könnte, dass $\forall y\; (\forall \;F(x,y) \wedge \forall x \;S(x))$, was bedeuten würde, dass jeder Mensch Student ist. 
\item $F_3$ bedeutet nicht „Wer ein Fach studiert, ist ein Student“. Es sagt vielmehr „Es existiert ein ein $y$ für jeden Menschen, das sagt, dass er nicht studiert“.
\end{itemize}
$\to$ in 2(a) sind wahre Voraussetzungen angenommen: Dabei sind $F_1$ und $F_3$ korrekter Weise wahr.\\
in 2(b) sind falsche Voraussetzungen angenommen: Dabei ist $F_1$ korrekter Weise falsch und $F_3$ fälschlicher Weise wahr.

\section{Aufgabe 3}
\begin{anumerate}
\item Jeder Informatikstudent studiert Informatik.\\
$\forall x \; (I(x) \to F(x,inf))$
\item Jeder Informatikstudent ist ein Student, aber nicht jeder Student ist ein Informatikstudent.\\
$\forall x \; \exists y \; ((I(x) \to S(x)) \wedge S(y) \wedge \neg I(y))  $\\
Achtung: $\forall x \; \exists y \;((I(x) \to S(x)) \wedge (S(y) \to \neg I(y)))$ geht nicht, da man nun für $y$ „Stuhl“ einsetzen könnte und die Formel richtig wäre da die Prämisse der zweiten Teilformel dann falsch und damit die Konklusion wahr wäre.
\item Jeder Student studiert ein Fach.\\
$\forall x \; \exists y \; (S(x) \to F(x,y))$
\item Wer ein Fach studiert, ist ein Student.\\
$\forall x \; \forall y \; (F(x,y) \to S(x))$\\
Achtung: $\forall x \; \exists y \; (F(x,y) \to S(x))$ geht nicht, da man nun für $y$ irgendetwas eingeben und die Prämisse wäre falsch, also die Konklusion richtig -- was aber falsch ist.
\item Es gibt kein Fach, das kein Student studiert.\\
$\forall y \; \exists x \; (F(x,y) \wedge S(x))$\\
Achtung: $\forall y \; \exists x \; (F(x,y) \to S(x))$ geht nicht, da man nun für $x$ irgendetwas einsetzen kann und dafür bei einem Fach, dem kein Student zugeordnet ist, die Prämisse falsch ist und somit die Aussage fälschlicher Weise wahr.
\end{anumerate}

\section{Aufgabe 4}
\begin{anumerate}
\item Es gibt ein Haustier, das fliegen kann.\\
$\exists x\; (haustier(x) \wedge fliegt (x))$
\item Katzen können nicht fliegen.\\
$\forall x\; (katze(x) \to \neg fliegt(x))$
\item Wenn ein Haustier fliegen kann, dann ist es ein Vogel.\\
$\forall x\; (haustier(x) \wedge fliegt(x) \to vogel(x))$
\item  Kein Haustier ist sowohl Katze als auch Vogel.\\
$\forall x\; (haustier(x) \to \neg (katze(x) \wedge vogel(x)))$
\end{anumerate}

\chapter{Praktikum 2: Prädikatenlogik, Prolog}
\section{Aufgabe 5}
\begin{align*}
&\forall x \;\exists y \;P(x,y) &\wedge&\\
&\forall x \; \forall y \; (P(x,y) \to \neg P(y,x)) &\wedge&\\
&\forall x \; \forall y \; \forall z \; (P(x,y) \wedge P(y,z) \to P(x,z))&&
\end{align*}
\begin{itemize}
\item $P()$ ist nicht kommutativ
\item $P()$ ist transitiv
\end{itemize}
$P(x,y)$: $x$ ist kleiner als $y$.\\
$P=\{(x,y) \;|\; x< y\}$\\
$\Rightarrow$ $x=1,\; y=2,\; z=3$

\section{Aufgabe 6}
$x$: Straße, $y$: Ort\\
$\forall x \; \forall y \; ( regnet(x) \wedge strasseinOrt(y,x) \to nass(y) )$\bigskip\\
$ort(dresden) \wedge regnet(dresden) $\\
$strasseInOrt(hochschulstrasse,dresden)$

\begin{lstlisting}[language=Prolog]
ort(dresden).	% nicht nötig
strasse(hochschulstrasse).	% nicht nötig
strasseInOrt(hochschulstrasse,dresden).
regnet(dresden).

nass(Y) :- regnet(X), strasseInOrt(Y,X).
\end{lstlisting}

\section{Aufgabe 7}
Bedeutungen:\\
\begin{tabular}{l l}
$terrorv(x)$ & $x$ ist eine terroristische Vereinigung\\
$terror(x)$ & $x$ ist ein Terrorist\\
$org(x)$ & $x$ ist eine Organisation\\
$mitgl(x,y)$ & $x$ ist Mitglied in Organisation $y$\\
$anf(x,y)$ & $x$ ist Anführer der Organisation $y$\\
$verkauf(x,y,z)$ & $x$ hat $y$ an $z$ verkauft\\
$waffe(x)$ & $x$ ist eine Waffe\\
$strafbar(x)$ & $x$ macht sich strafbar\\
$auto(x)$ & $x$ ist ein Auto
\end{tabular}
\paragraph*{Teil a}
\begin{itemize}
\item $\forall x\; \forall y\;(anf(x,y) \wedge terror(x) \to terrorv(y) )$
\item $\forall x\;\forall y\; (mitgl(x,y) \wedge terrorv(y) \to terror(x))$
\item $\exists x\; (anf(x,\text{IS}) \wedge terror(x))$
\item $mitgl(\text{Abul-Hasan Al Muhajir},\text{IS})$
\item $\forall x\; \forall y\; \forall z ( verkauf(x,y,z) \wedge terror(z) \wedge waffe(y) \to strafbar(x))$
\item $verkauf(\text{Omar Abdul Al-Hassan},\text{M16},\text{Abul-Hasan Al Muhajir})$
\item $waffe(\text{M16})$
\item $mitgl(\text{August Meier})$
\item $org(\text{ADAC})$
\item $verkauf(\text{Bernd Müller}, \text{Auto},\text{August Meier})$
\end{itemize}

\paragraph*{Teil b}
Prolog-Programm:
\begin{lstlisting}[language=Prolog]
terrorv(Y) :- anf(X,Y), terror(X).
terror(X) :- mitgl(X,Y) , terrorv(Y).
terror("Terrorist").
org("ADAC").
mitgl("August Meier", "ADAC").
mitgl("Abul-Hasan Al Muhajir", "IS").
anf("Terrorist","IS").
strafbar(X) :- verkauf(X,Y,Z), waffe(Y), terror(Z).
verkauf("Omar Abdul Al-Hassan", "M16", "Abul-Hasan Al Muhajir").
verkauf("Bernd Müller", "Auto", "August Meier").
waffe("M16").
\end{lstlisting}
\begin{enumerate}[label=(\roman*)]
\item Abfrage, ob Omar Abdul Al-Hassan strafbar:
\begin{lstlisting}[language=Prolog]
?- strafbar("Omar Abdul Al-Hassan").
true.
\end{lstlisting}
\item Abfrage, ob Berdn Müller strafbar:
\begin{lstlisting}[language=Prolog]
?- strafbar("Bernd Müller").
false.
\end{lstlisting}
\end{enumerate}
\paragraph*{Teil c}
Suchbaum:\\
\resizebox{.99\textwidth}{!}{
\begin{tikzpicture}[scale=1.1]
\node (v1) at (0,0) {\lstinline[language=Prolog]|strafbar("Omar Abul Al-Hassan")|};
\node (v2) at (-5,-2) {\lstinline[language=Prolog]|verkauf("Omar Abdul Al-Hassan", Y, Z)|};
\node (v4) at (0,-2) {\lstinline[language=Prolog]|waffe(Y)|};
\node (v6) at (4,-2) {\lstinline[language=Prolog]|terror(Z)|};
\node (v3) at (-5,-4) {\lstinline[language=Prolog]|verkauf ("Omar Abdul Al-Hassan", "M16", "Abul-Hasan Al Muhajir")|} (v3) node[below]{\lstinline[language=Prolog]|true|};
\draw (v1) -- (v2) -- (v3);
\node (v5) at (0,-6) {\lstinline[language=Prolog]|waffe("M16")|} (v5) node[below]{\lstinline[language=Prolog]|true|};
\draw (v1) -- (v4) -- (v5);
\node (v7) at (4,-6) {\lstinline[language=Prolog]|terror("Abul-Hasan Al Muajir")|};
\draw (v1) -- (v6) -- (v7);
\node (v8) at (2,-8) {\lstinline[language=Prolog]|mitgl("Abul-Hasan Al Muajir",Y)|};
\node (v10) at (6,-8) {\lstinline[language=Prolog]|terrorv(Y)|};
\node (v9) at (2,-10) {\lstinline[language=Prolog]|mitgl("Abul-Hasan Al Muhajir", "IS")|} (v9) node[below]{\lstinline[language=Prolog]|true|};
\draw (v7) -- (v8) -- (v9);
\node (v11) at (6,-12) {\lstinline[language=Prolog]|terrorv("IS")|};
\node (v12) at (4,-14) {\lstinline[language=Prolog]|anf(X,"IS")|};
\node (v14) at (8,-14) {\lstinline[language=Prolog]|terror(X)|};
\node (v13) at (4,-16) {\lstinline[language=Prolog]|anf("Terrorist","IS")|}(v13) node[below]{\lstinline[language=Prolog]|true|};
\node (v15) at (8,-18) {\lstinline[language=Prolog]|terror("Terrorist")|}(v15) node[below]{\lstinline[language=Prolog]|true|};
\draw (v7) -- (v10) -- (v11) -- (v12) -- (v13);
\draw (v11) -- (v14) -- (v15);
\def\xa{-26};
\def\ra{1cm};
\draw[shift=(\xa : \ra)] (0,0) arc (\xa:-158:\ra);
\def\xb{-45};
\def\rb{1cm};
\draw[shift=(\xb : \rb)] (4,-6) arc (\xb:-135:\rb);
\def\xc{-45};
\def\rc{1cm};
\draw[shift=(\xc : \rc)] (6,-12) arc (\xc:-135:\rc);
\end{tikzpicture}
}

\chapter{Praktikum 3: Unifizierung, Prolog}

\section{Aufgabe 8}
$a,f(a)$ … Konstanten
$x,y$ … Variablen
\begin{anumerate}
\item $P(x,y,a),\; P(x,x,x)$ sind unifizierbar durch $x/a,\; y/a$
\item $P(f(a),x,y),\; P(x,a,a)$ sind nicht unifizierbar, da $x$ sowohl mit $f(a)$ und $a$ ersetzt werden müsste (ersten beiden Positionen)
\item $P(f(a),a,x),\; P(x,x,a)$ sind nicht unifizierbar,  da $x$ sowohl mit $f(a)$ und $a$ ersetzt werden müsste (ersten beiden Positionen)
\end{anumerate}

\section{Aufgabe 9}
\lstinputlisting[language=Prolog]{Prolog/aufg9.pl}
Stammbaum:
\begin{center}
\begin{tikzpicture}
\node (v1) at (-9.5,-6.5) {Tim};
\node (v2) at (-10.5,-5) {Jens};
\node (v3) at (-8.5,-5) {Heike};
\node (v4) at (-12,-5) {Jutta};
\node (v5) at (-13.5,-5) {Julia};
\node (v6) at (-7,-5) {Jan};
\node (v7) at (-5,-5) {Eva};
\node (v8) at (-7,-6.5) {Anne};
\node (v9) at (-5,-6.5) {Anke};
\node (v10) at (-11,-3.5) {Robert};
\node (v11) at (-9,-3.5) {Sandra};
\node (v12) at (-14,-2) {Wilhelm};
\node (v13) at (-12,-2) {Frieda};
\node (v14) at (-10,-2) {Kurt};
\node (v15) at (-8,-2) {Lisa};
\node (v16) at (-4.5,-3.5) {Thomas};
\node (v17) at (-1.5,-3.5) {Christine};
\node (v18) at (-6,-2) {Winfried};
\node (v19) at (-3,-2) {Elke};
\node (v20) at (-1.5,-2) {Hans};
\node (v21) at (0.5,-2) {Paula};
\node (v22) at (-13,-2) {};
\node (v24) at (-9,-2) {};
\node (v23) at (-10,-3.5) {};
\node (v25) at (-9.5,-5) {};
\node (v26) at (-6,-5) {};
\node (v27) at (-3,-3.5) {};
\node (v28) at (-4.5,-2) {};
\node (v29) at (-0.5,-2) {};
\draw (v12) -- (v13);
\draw (v14) -- (v15);
\draw (v18) -- (v19);
\draw (v20) -- (v21);
\draw (v16) -- (v17);
\draw (v10) -- (v11);
\draw (v2) -- (v3);
\draw (v6) -- (v7);
\draw (v5) -- (v23);
\draw (v4) -- (v23);
\draw (v2) -- (v23);
\draw (v1) -- (v25);
\draw (v8) -- (v26);
\draw (v9) -- (v26);
\draw (v6) -- (v23);
\draw (v7) -- (v27);
\draw (v16) -- (v28);
\draw (v17) -- (v29);
\draw (v10) -- (v22);
\draw (v11) -- (v24);
\end{tikzpicture}
\end{center}
Beispiel-Abfragen:
\begin{lstlisting}[language=Prolog]
?- bruder(X, julia).
X = jan ;
X = jens .

?- schwester(X,jens).
X = julia ;
X = jutta .

?- bruder(jens, X).
X = jan ;
X = julia ;
X = jutta .

?- enkel(X,robert).
X = anne ;
X = anke ;
X = tim .

?- tante(X,tim).
X = julia ;
X = jutta.

?- tante(X,tim).
X = julia ;
X = jutta . % es fehlt eva?!
\end{lstlisting}

\chapter{Praktikum 4: Transitive Hülle, Prolog}

\section{Aufgabe 11}
Der Suchbaum würde sich in einer Rekursion aufhängen: beim erste \lstinline$weg()$.\\
Lösung: Eine extra Funktion für die transitive Hülle einführen, die die Rekursion auf den rechten Zweig verschiebt, wo sie sich auflösen oder zu einem Fehler führen kann.
\begin{lstlisting}[language=Prolog]
tweg(X,Y) :- weg(X,Y); weg(X,Z), tweg(Z,Y).
\end{lstlisting}

\section{Aufgabe 12}
???

\chapter{Praktikum 5: Vereinfachung Prolog}

\section{Aufgabe 13}
\lstinputlisting[language=Prolog]{Prolog/aufg13.pl}















\end{document}